\section{Структура работы}

\hfill%
\begin{minipage}{0.5\linewidth}
	\textit{«Реши, что ты хочешь: защититься или народ удивить».}
	
	\hfill --- народная мудрость
\end{minipage}

\subsection{Общие положения}

Выпускная квалификационная работа бакалавра и магистранта (а если брать шире, то и аспиранта) имеет классическую четыёхчастную структуру: анализ проблемы и/или анализ традиционного процесса решения задачи, разработка нового способа решения задачи, выбор и обоснование выбора средств реализации решения, реализация решения.
Так как  ВКР посвящена разработке, в первую очередь, программного обеспечения, то удобно для организации её структуры использовать промышленный способ создания автоматизированных систем (АС) (рис. \ref{fig:AS_MAIT}), описываемый в работах, посвящённых методологии автоматизации интеллектуального труда (МАИТ). \cite{gd:nm}.   

\begin{figure}
	%\resizebox{\linewidth}{!}{
	\centering
	\begin{tikzpicture}
		[
		  pd/.style={draw, align=center, minimum height = 1cm},
		  pdnumber/.style={above right, yshift=-3pt, xshift=-1pt},
		  pdrole/.style={below right, yshift=2pt, xshift=-1pt},
		  node distance=2em and 2.2em,
		]
				
		\fontsize{10pt}{9pt}\selectfont
		
		\draw node[pd] (z11) {Промышленный способ\\создания АС};
		
		\draw node[pd, below = of z11, xshift=-5cm, fill=green!30!white] (z21) {Предпроектное\\обследование};
		\draw node[pd, right = of z21] (z22) {Проектирование АС};
		\draw node[pd, right = of z22] (z23) {Подготовка \\реализации АС};
		\draw node[pd, right = of z23] (z24) {Реализация АС};
		\draw node[pd, below = of z21, xshift=-0.3cm, fill=red!30!white] (z31) {Концептуальное\\моделирование};
		\draw node[pd, right = of z31, fill=orange!30!white ] (z32) {Инфологическое\\моделирование}; 
		\draw node[pd, right = of z32 ] (z33) {Выбор программно-технических\\средств реализации}; 
		\draw node[pd, right = of z33, fill=Magenta!30!white] (z34) {Даталогическое\\моделирование}; 
		
		\draw [-Latex] ($(z11.south west)!0.5!(z11.south)$) -- ++ (down:1em) -| (z21.90);
		\draw [Latex-] ($(z11.south east)!0.5!(z11.south)$) -- ++ (down:1em) -| (z24.90);
		
		\draw [-Latex] (z21) -- (z22);
		\draw [-Latex] (z22) -- (z23);
		\draw [-Latex] (z23) -- (z24);
		
		\draw [-Latex] ($(z22.south west)!0.5!(z22.south)$) -- ++ (down:1em) -| (z31.90);
		\draw [Latex-] ($(z22.south east)!0.5!(z22.south)$) -- ++ (down:1em) -| (z32.70);
		
		\draw [-Latex] (z31) -- (z32);
		
		\draw [-Latex] ($(z23.south west)!0.5!(z23.south)$) -- ++ (down:1em) -| (z33.120);
		\draw [Latex-] ($(z23.south east)!0.5!(z23.south)$) -- ++ (down:1em) -| (z34.90);
		
		\draw [-Latex] (z33) -- (z34);  
		
		\foreach \z in {z11,z21,z22,z23,z24,z31,z32,z33,z34}
		{
			\draw (\z.south east) node[pdnumber] {\z};
		}
		
		\foreach \z in {z11,z22,z23}
		{
			\draw (\z.north east) node [pdrole]{[};
		}
		
		\foreach \z in {z24,z32,z34}
		{
			\draw (\z.north east) node [pdrole]{]};
		}	
		
	\end{tikzpicture} %}
	\caption[х]{Промышленный способ создания в соответствии с МАИТ \footnotemark}
	\label{fig:AS_MAIT}
\end{figure} 
\footnotetext{Вот тема на ВКР --- разработка библиотеки для \LaTeX{} и Ti\textit{k}Z для рисования подобных схем. Жду предложений.}

В соответствии с МАИТ разработка автоматизированной системы происходит в 4 укрупненных этапа (а в ВКР 4 главы, совпадение?):
\begin{enumerate}
	\item Предпроектное обследование.
	
	\item Проектирование.
	
	\item Подготовка реализации.
	
	\item Реализация.
	
\end{enumerate}  

\subsection{Цель, задачи, объект и предмет}

\subsubsection{Формулирование цели работы}

Цель --- это положительный эффект который должен быть достигнут в результате выполнения работы. Хорошо сформулированная цель должна показывать какая реальная характеристика или какой показатель в какую сторону и на сколько и за счёт чего изменится. Самая частая  ошибка писать в цели: <<разработка приложения~...>> или <<создание информационной системы~...>> и т.д. Это не цель, а средство ее достижения.

Примеры хорошо сформулированных целей:
\begin{itemize}
	\item снижение трудозатрат сотрудника проектно-конструкторского бюро за счёт автоматизации процесса анализа данных, поступающих с производственного оборудования;
	
	\item повышение коэффициента покрытия кода автотестами за счёт  перехода на разработку через написание предварительных тестов (Test Driven Development, TDD);
	
	\item cнижение ежемесячных расходов на облачную инфраструктуру за счёт перехода с серверов общего назначения на специализированные ARM-серверы.
\end{itemize}

Можно сформулировать несколько условных шаблонов подобных целей. Так как  целеполагание под каждую задачу обладает своими особенности шаблоны носят условный характер.
\begin{enumerate}
	\item <<Снижение/Рост [Показатель] у [Субъект/Отдел] на [X\% или до Y значения] за счёт [Конкретное действие/Инструмент] >>.
	\item <<Ускорение/Замедление [Процесс][Какого?]  [или для Кого или чего] за счёт [Автоматизация/Изменение алгоритма/Разработка]>>.
	\item <<Повышение/Снижение [Качество/Количество] [Продукта/Услуги] для [Клиент/Сотрудник] за счёт [Новая методика/Новая система/программа ...]>>.
\end{enumerate} 

В конце работы необходимо будет показать как заявленная цель была достигнута --- привести количественные и качественные показатели: <<было так, стало так>>.

Отдельно стоит рассмотреть цели вида <<повышение эффективности [процесса, субъекта, ... ]>>. Такая формулировка обязывает автора пояснить, что такое <<эффективность>>, в каких единицах она измеряется и придумать способ в конце работы продемонстрировать, как она была повышена.

Не стоит использовать в цели такие характеристики, которые нельзя показать без цифр: улучшить, оптимизировать\footnotemark, усилить, активизировать, интенсифицировать, повысить эффективность (без цифр), сделать лучше, навести порядок. Это неизменяемые показатели, которые вообще лучше избегать в технической литературе. 

\footnotetext{Термин <<оптимизация>> можно использовать, если имеет место именно математическая оптимизация~--- поиск экстремума целевой функции, применение методов оптимизации, поиск критериев оптимизации. } 

\subsubsection{Объект и предмет}

Объект исследования --- это область, явление, сфера знаний, процесс, в рамках которых будет осуществляться исследование.

Предмет исследования –-- более детализированное и узкое понятие, которое обязательно должно быть частью объекта и не может выходить за его рамки. 


\begin{table}
	\caption{Примеры объектов  и предметов}%
	\label{tbl:objectsubject}
	\begin{tabularx}{\linewidth}{|>{\raggedright}X|>{\raggedright}X|>{\raggedright\arraybackslash}X|}
	\hline
	\centering \textbf{Тема}	& \centering \textbf{Объект}  & \centering \textbf{Предмет} \arraybackslash \\ \hline
	Исследование и анализ графических нотаций методов моделирования и разработка модуля визуализации 3D-диаграмм для интегрированной среды «ИС-2» &
	Графические нотации для визуализации статических и динамических структур различных методов моделирования &
	Алгоритмы и технологии визуализации трёхмерной диаграммы, отражающей взаимосвязь динамических и статических структур концептуальных моделей 
	\\ \hline
	Исследование особенностей организации и формирования технологической документации для выделения её содержательных и формальных характеристик &
	Технологическая документация предприятий &
	Cодержательные и формальные характеристики технологической документации предприятий 
	\\ \hline
	Исследование процесса формирования структуры классов на
	основе инфологической модели и разработка средств его поддержки &
	Даталогическое моделирование (в рамках объектно-
	ориентированного подхода)&
	Процедура отображения инфологической модели в даталогическую
	\\ \hline
	Разработка автоматизированной подсистемы
	расчёта стыковых сварных соединений машиностроительных
	конструкций &
	Процесс проектирования стыковых сварных соединений, выполняемых дуговыми методами сварки&
	Алгоритмы и программные средства автоматизации расчёта геометрических параметров стыковых сварных соединений 
	\\ \hline
	
	\end{tabularx}
\end{table}


\subsection{Описание глав ВКР}

\subsection{Первая глава ВКР}

\subsection{Вторая глава ВКР}

\subsection{Третья глава ВКР}

\subsection{Четвёртая глава ВКР}